\section{Exercise 1: Conditional Expectation}

\textbf{Problem Statement:} Let $X$ be a random variable taking values in some finite set and $Y \in L^1(\Omega, \mathcal{F}, P)$. Prove that there exist $n \in \mathbb{N}$ and a partition $(A_1, \dots, A_n)$ of $\Omega$ such that:
\[ E[Y|X] = \sum_{i=1}^n \frac{E[Y 1_{A_i}]}{P[A_i]} 1_{A_i} \quad P\text{-a.s.} \]

\subsection*{Step 1: Constructing the Partition}
We are told that $X$ takes values in a finite set. Let's call these distinct values $\{a_1, a_2, \dots, a_n\}$. By definition, this means that for any $i \neq j$, $a_i \neq a_j$.

We can group the outcomes in our sample space $\Omega$ based on the value $X$ takes. Let's define the set $A_i$ as the event where $X$ equals $a_i$:
\[ A_i = X^{-1}(a_i) = \{\omega \in \Omega \mid X(\omega) = a_i\} \]
for $i \in \{1, \dots, n\}$.

Because $X$ must take exactly one of these values for any given outcome $\omega$, the sets $A_1, \dots, A_n$ do not overlap (they are mutually exclusive) and together they cover the entire sample space $\Omega$. Therefore, $(A_1, \dots, A_n)$ forms a valid partition of $\Omega$. 

We can also write $X$ using these partition sets:
\[ X = \sum_{i=1}^n a_i 1_{A_i} \]

\subsection*{Step 2: Checking the Measurability Condition}
To prove that our proposed formula is indeed the conditional expectation $E[Y|X]$, it must satisfy two mathematical conditions. The first is that the proposed random variable must be $\sigma(X)$-measurable.

The $\sigma$-algebra generated by $X$, denoted $\sigma(X)$, is exactly the $\sigma$-algebra generated by our partition: $\sigma(X) = \sigma(A_1, \dots, A_n)$. 

Let's look at our proposed formula:
\[ Z = \sum_{i=1}^n \frac{E[Y 1_{A_i}]}{P[A_i]} 1_{A_i} \]
Notice that $\frac{E[Y 1_{A_i}]}{P[A_i]}$ is simply a constant (a real number) for each $i$. Since the indicator functions $1_{A_i}$ are $\sigma(X)$-measurable (because $A_i \in \sigma(X)$), any linear combination of them is also $\sigma(X)$-measurable. Thus, the first condition is satisfied.

\subsection*{Step 3: Checking the Partial Expectation Condition}
The second condition for conditional expectation requires that for any set $A \in \sigma(X)$, the expectation of our proposed formula over the set $A$ equals the expectation of $Y$ over the set $A$. Mathematically, for all $A \in \sigma(X)$:
\[ E[Z 1_A] = E[Y 1_A] \]

Since $\sigma(X)$ is generated by the partition $(A_1, \dots, A_n)$, it is sufficient to check this condition just for the building blocks of the $\sigma$-algebra: the partition sets $A_j$ for $j \in \{1, \dots, n\}$.

Let's plug our formula $Z$ into the left side of the equation:
\[ E[Z 1_{A_j}] = E\left[ \left( \sum_{i=1}^n \frac{E[Y 1_{A_i}]}{P[A_i]} 1_{A_i} \right) 1_{A_j} \right] \]

By the linearity of the expectation operator, we can move the expectation inside the sum:
\[ = \sum_{i=1}^n \frac{E[Y 1_{A_i}]}{P[A_i]} E[1_{A_i} 1_{A_j}] \]

Now, remember that our sets $A_i$ and $A_j$ are disjoint if $i \neq j$. Therefore, the product of their indicator functions $1_{A_i} 1_{A_j} = 0$ everywhere except when $i = j$. When $i = j$, $1_{A_j} 1_{A_j} = 1_{A_j}$. This makes the entire sum collapse into a single term where $i = j$:
\[ = \frac{E[Y 1_{A_j}]}{P[A_j]} E[1_{A_j}] \]

Since $E[1_{A_j}]$ is simply the probability of the event $A_j$, we have $E[1_{A_j}] = P[A_j]$. Substituting this in gives:
\[ = \frac{E[Y 1_{A_j}]}{P[A_j]} P[A_j] = E[Y 1_{A_j}] \]

We have successfully shown that $E[Z 1_{A_j}] = E[Y 1_{A_j}]$, which satisfies the second condition. Therefore, our formula $Z$ is indeed $E[Y|X]$.